\documentclass[12pt,a4paper]{article}

% ============================================================
% PACKAGES
% ============================================================
\usepackage[margin=2.5cm]{geometry}
\usepackage{amsmath,amssymb}
\usepackage{booktabs}
\usepackage{array}
\usepackage{longtable}
\usepackage{graphicx}
\usepackage{float}
\usepackage{fancyhdr}
\usepackage{titlesec}
\usepackage{tocloft}
\usepackage{hyperref}
\usepackage{listings}
\usepackage{xcolor}
\usepackage{parskip}
\usepackage{enumitem}
\usepackage{microtype}
\usepackage{caption}
\usepackage{subcaption}
\usepackage{multirow}
\usepackage{makecell}

% ============================================================
% HYPERREF SETUP
% ============================================================
\hypersetup{
    colorlinks=true,
    linkcolor=black,
    filecolor=black,
    urlcolor=blue,
    pdftitle={CHG 433 Assignment II -- Complete Report},
    pdfauthor={Group 9},
}

% ============================================================
% CODE LISTING STYLE
% ============================================================
\lstdefinestyle{matlabstyle}{
    language=Matlab,
    basicstyle=\ttfamily\footnotesize,
    keywordstyle=\color{blue}\bfseries,
    commentstyle=\color{green!50!black},
    stringstyle=\color{red!70!black},
    numbers=left,
    numberstyle=\tiny\color{gray},
    stepnumber=1,
    numbersep=8pt,
    backgroundcolor=\color{gray!5},
    showspaces=false,
    showstringspaces=false,
    breaklines=true,
    frame=single,
    rulecolor=\color{gray!40},
    tabsize=4,
    captionpos=b,
}

% ============================================================
% PAGE STYLE
% ============================================================
\pagestyle{fancy}
\fancyhf{}
\rhead{\small CHG 433 -- Chemical Engineering Analysis}
\lhead{\small Group 9 -- Assignment II}
\rfoot{\thepage}
\renewcommand{\headrulewidth}{0.4pt}

% ============================================================
% SECTION FORMATTING
% ============================================================
\titleformat{\section}[block]{\large\bfseries\centering}{}{0em}{}[\vspace{0.5ex}\hrule\vspace{1ex}]
\titleformat{\subsection}[block]{\normalsize\bfseries}{}{0em}{}
\titleformat{\subsubsection}[block]{\normalsize\itshape\bfseries}{}{0em}{}

% ============================================================
% DOCUMENT BEGIN
% ============================================================
\begin{document}

% ============================================================
% TITLE PAGE
% ============================================================
\begin{titlepage}
  \centering
  \vspace*{1cm}

  {\LARGE \textbf{CHEMICAL ENGINEERING ANALYSIS}}\\[0.4cm]
  {\Large Assignment II}\\[1.5cm]

  {\large Department of Chemical Engineering\\
  University of Lagos, Akoka, Lagos\\[0.5cm]
  CHG 433: Chemical Engineering Analysis\\[0.5cm]
  2025/2026 Academic Session\\[1.5cm]}

  {\large \textbf{Group 9}}\\[2cm]

  \begin{tabular}{ll}
    \toprule
    \textbf{Member Name} & \textbf{Matric Number} \\
    \midrule
    Group 9 Member 1 & \\
    Group 9 Member 2 & \\
    Group 9 Member 3 & \\
    Group 9 Member 4 & \\
    \bottomrule
  \end{tabular}

  \vfill

  {\large \today}
\end{titlepage}

% ============================================================
% ATTESTATION PAGE
% ============================================================
\newpage
\thispagestyle{plain}
\begin{center}
  {\large \textbf{ATTESTATION}}
\end{center}

The authors of this paper are the members of Group 9.\\[1cm]

\noindent\begin{tabular}{p{6cm}p{4cm}p{4cm}}
Group 9 Member 1 & \underline{\hspace{4cm}} & \\[0.8cm]
Group 9 Member 2 & \underline{\hspace{4cm}} & \\[0.8cm]
Group 9 Member 3 & \underline{\hspace{4cm}} & \\[0.8cm]
Group 9 Member 4 & \underline{\hspace{4cm}} & \\[0.8cm]
\end{tabular}

\vspace{1cm}
\noindent By the signatures above, each author attests to the fact that the contents of this paper are,
to the best of our knowledge, neither plagiarised nor false.

\vspace{2cm}
\noindent \textbf{Contributions}

\vspace{0.5cm}
\noindent\begin{tabular}{p{4cm}p{10cm}}
\toprule
\textbf{Member} & \textbf{Roles / Contributions} \\
\midrule
Group 9 Member 1 & Layout/typing, results analysis and interpretation \\[0.4cm]
Group 9 Member 2 & Programming/computing tasks, interpretation, discussion \\[0.4cm]
Group 9 Member 3 & Conceptual formulation of solutions, programming \\[0.4cm]
Group 9 Member 4 & Conceptual formulation of solutions, programming \\
\bottomrule
\end{tabular}

% ============================================================
% EXECUTIVE SUMMARY
% ============================================================
\newpage
\begin{center}
  {\large \textbf{EXECUTIVE SUMMARY}}
\end{center}

Mathematical modelling and analysis are of fundamental importance in Chemical Engineering.
Engineers must design specialised reactors and processes to create new substances or materials
while adhering to strict safety regulations and working to maximise efficiency from both
economic and ecological standpoints. The aim of this assignment is to establish a firm
knowledge in analysing, modelling and refining new and existing systems, establishing the
balance between complexity and accuracy, and to foster teamwork between members towards
creating or advancing industrial processes.

Six questions were solved and evaluated using numerical methods and the computer application
MATLAB. The solutions were first computed manually by making assumptions and applying
theories; then the required formulas and lines of code were typed into MATLAB for
verification.

\textbf{Question 1} considered a three-reactor system operating at steady state. By making a mass
balance around each reactor, a system of three simultaneous linear equations was generated
and solved for concentrations $x_A$, $x_B$, and $x_C$ using Cramer's Rule, the Matrix Inverse
Method, and the MATLAB backslash operator. The concentrations converged to
$x_A = 14.67\,\text{mg/m}^3$, $x_B = 7.33\,\text{mg/m}^3$, and $x_C = 10.10\,\text{mg/m}^3$.

\textbf{Question 2} was in two parts. Part A determined the temperature in a Resistance Temperature
Detector (RTD) when its resistance was 300\,$\Omega$, using the Newton-Raphson and Interval
Bisection methods, yielding a temperature of approximately $260.22\,^{\circ}\text{C}$. Part B determined
the volume of a real gas (N\textsubscript{2}) at multiple pressures using the Van der Waals equation, with
MATLAB's \texttt{fzero} function confirming the inverse pressure-volume relationship.

\textbf{Question 3} derived a model relating inlet volumetric flow to liquid level in a tank, identified
it as a lumped-parameter first-order ODE, found the initial steady-state level and derived a
closed-form dynamic response. The ODE was then solved numerically using the Explicit
Euler and Runge-Kutta 4\textsuperscript{th}-Order (RK4) methods for step lengths of 0.01, 0.05, 0.1,
0.5, 1, and 10\,s.

\textbf{Question 4} applied the method of least squares to fit a function of the form
$\ln(k) = C + b\ln(T) - D/T$ to experimental rate-coefficient data. The fitted constants were
used to deduce the Arrhenius parameters $A$ and $E_a$ for the reaction $\text{CH}_4 + \text{O} \to \text{CH}_3 + \text{OH}$.

\textbf{Question 5} solved a flash vaporization problem for a ternary mixture of n-butane,
n-pentane, and n-hexane at 10\,atm and 270\,$^{\circ}$F using the Rachford-Rice equation and
Wilson's K-value correlation, determining the vapour fraction $\beta$, vapour flow $V$, liquid
flow $L$, and phase compositions $x_i$ and $y_i$.

\textbf{Question 6} determined the cost per pound of five trail-mix ingredients by formulating and
solving a $5\times5$ linear system from five known mixture-cost equations using the MATLAB
backslash operator and back-substitution for validation.

% ============================================================
% TABLE OF CONTENTS
% ============================================================
\newpage
\tableofcontents
\newpage

% ============================================================
% INTRODUCTION
% ============================================================
\section*{INTRODUCTION}
\addcontentsline{toc}{section}{Introduction}

Mathematical techniques and models are used in almost every element of a chemical
engineer's professional life. Modern industrial production is based on solid scientific methods,
many of which are part of chemical engineering. Engineers must design specialised reactors and
processes to create new substances or materials while adhering to strict safety regulations and
working to maximise efficiency from both an economic and ecological standpoint.

According to the American Institute of Chemical Engineers (AIChE), mathematical modelling
is a powerful tool for understanding, designing, and predicting processes and process
equipment in the chemical industry, including the conservation of momentum, energy, and
material. Models are designed to predict and optimise the behaviour of engineering operations.

Modern chemical engineers are not only expected to be adept at using these techniques; they
are also expected to be able to apply computer-based solutions when analytical ones are
impractical. This assignment demonstrates both approaches—manual analytical methods and
MATLAB computational implementations—across six engineering problems drawn from
linear algebra, root-finding, ordinary differential equations, regression, thermodynamic
flash calculations, and simultaneous linear systems.

\newpage

% ============================================================
% SECTION 1: QUESTION 1
% ============================================================
\section{SECTION 1 -- Question 1: Three-Reactor Steady-State System}

\textbf{Problem Statement:} A steady-state mass balance on a three-reactor system produces the
following system of linear equations for concentrations $x_A$, $x_B$, and $x_C$ (in mg/m\textsuperscript{3}):

\begin{align}
  120x_A - 60x_B + 0x_C &= 1320 \label{eq:q1a}\\
  40x_A - 80x_B + 0x_C &= 0 \label{eq:q1b}\\
  80x_A + 20x_B - 150x_C &= -195 \label{eq:q1c}
\end{align}

\textbf{Required:} Solve the system for $x_A$, $x_B$, and $x_C$.

% ----------------------------------------------------------
\subsection{1.1 Manual Calculation}

The system \eqref{eq:q1a}--\eqref{eq:q1c} can be expressed in matrix form as
$\mathbf{A}\mathbf{x} = \mathbf{b}$, where

\[
\mathbf{A} = \begin{bmatrix} 120 & -60 & 0 \\ 40 & -80 & 0 \\ 80 & 20 & -150 \end{bmatrix},
\quad
\mathbf{x} = \begin{bmatrix} x_A \\ x_B \\ x_C \end{bmatrix},
\quad
\mathbf{b} = \begin{bmatrix} 1320 \\ 0 \\ -195 \end{bmatrix}.
\]

% ---- 1.1.1 Cramer's Rule ----
\subsubsection{1.1.1 Cramer's Rule}

By Cramer's Rule:
\[
x_A = \frac{D_{x_A}}{D}, \quad x_B = \frac{D_{x_B}}{D}, \quad x_C = \frac{D_{x_C}}{D},
\]
where $D = \det(\mathbf{A})$.

\textbf{Determinant of A:}
\begin{align*}
D = \det(\mathbf{A}) &= 120[(-80)(-150) - (0)(20)] - (-60)[(40)(-150) - (0)(80)] + 0 \\
&= 120(12000) + 60(-6000) \\
&= 1\,440\,000 - 360\,000 = 1\,080\,000
\end{align*}

\textbf{Substituted determinants:}
\begin{align*}
D_{x_A} = \begin{vmatrix} 1320 & -60 & 0 \\ 0 & -80 & 0 \\ -195 & 20 & -150 \end{vmatrix}
&= 1320[(-80)(-150)-(0)(20)] - (-60)[(0)(-150)-(0)(-195)] \\
&= 1320(12000) + 0 = 15\,840\,000
\end{align*}

\begin{align*}
D_{x_B} = \begin{vmatrix} 120 & 1320 & 0 \\ 40 & 0 & 0 \\ 80 & -195 & -150 \end{vmatrix}
&= 120[(0)(-150)-(0)(-195)] - 1320[(40)(-150)-(0)(80)] \\
&= 0 - 1320(-6000) = 7\,920\,000
\end{align*}

\begin{align*}
D_{x_C} = \begin{vmatrix} 120 & -60 & 1320 \\ 40 & -80 & 0 \\ 80 & 20 & -195 \end{vmatrix}
&= 120[(-80)(-195)-(0)(20)] - (-60)[(40)(-195)-(0)(80)] + 1320[(40)(20)-(-80)(80)]\\
&= 120(15600) + 60(-7800) + 1320(800+6400)\\
&= 1\,872\,000 - 468\,000 + 9\,504\,000 = 10\,908\,000
\end{align*}

\textbf{Results:}
\[
\boxed{x_A = \frac{15\,840\,000}{1\,080\,000} = 14.67\,\text{mg/m}^3, \quad
x_B = \frac{7\,920\,000}{1\,080\,000} = 7.33\,\text{mg/m}^3, \quad
x_C = \frac{10\,908\,000}{1\,080\,000} \approx 10.10\,\text{mg/m}^3}
\]

% ---- 1.1.2 Matrix Inverse Method ----
\subsubsection{1.1.2 Matrix Inverse Method}

\[
\mathbf{x} = \mathbf{A}^{-1}\mathbf{b}
\]

The inverse of $\mathbf{A}$ is computed as $\mathbf{A}^{-1} = \frac{1}{\det(\mathbf{A})}\,\text{adj}(\mathbf{A})$.

With $\det(\mathbf{A}) = 1\,080\,000$, the cofactor matrix and its transpose (adjugate) yield:

\[
\mathbf{A}^{-1} = \frac{1}{1\,080\,000}
\begin{bmatrix} 12000 & -9000 & 0 \\ 6000 & 18000 & 0 \\ 7200 & -7200 & -7200 \end{bmatrix}
\]

\[
\mathbf{x} = \mathbf{A}^{-1}\mathbf{b} =
\frac{1}{1\,080\,000}
\begin{bmatrix} 12000 & -9000 & 0 \\ 6000 & 18000 & 0 \\ 7200 & -7200 & -7200 \end{bmatrix}
\begin{bmatrix} 1320 \\ 0 \\ -195 \end{bmatrix}
=
\begin{bmatrix} 14.67 \\ 7.33 \\ 10.10 \end{bmatrix}
\text{mg/m}^3
\]

% ---- 1.1.3 Gaussian Elimination ----
\subsubsection{1.1.3 Gaussian Elimination Method}

Starting augmented matrix:
\[
\left[\begin{array}{ccc|c}
120 & -60 & 0 & 1320\\
40 & -80 & 0 & 0\\
80 & 20 & -150 & -195
\end{array}\right]
\]

Applying row operation $R_1 \leftarrow R_1 - 3R_2$:
\[
\left[\begin{array}{ccc|c}
0 & 180 & 0 & 1320\\
40 & -80 & 0 & 0\\
80 & 20 & -150 & -195
\end{array}\right]
\]

Applying row operation $R_3 \leftarrow R_3 - 2R_2$:
\[
\left[\begin{array}{ccc|c}
0 & 180 & 0 & 1320\\
40 & -80 & 0 & 0\\
0 & 180 & -150 & -195
\end{array}\right]
\]

Back substitution from row 1 gives $x_B = 1320/180 = 7.33$ mg/m\textsuperscript{3}. From row 2: $x_A = (80 \times 7.33)/40 = 14.67$ mg/m\textsuperscript{3}. From row 3: $x_C = (-195 - 180\times7.33)/(-150) = 10.10$ mg/m\textsuperscript{3}.

\[
\boxed{x_A = 14.67\,\text{mg/m}^3,\quad x_B = 7.33\,\text{mg/m}^3,\quad x_C = 10.10\,\text{mg/m}^3}
\]

% ----------------------------------------------------------
\subsection{1.2 Software Calculation (MATLAB)}

The MATLAB script \texttt{Ass2025-2026-Grp9\_Q1.m} solves the system using three independent
approaches: the backslash operator (\texttt{A\textbackslash b}), matrix inverse (\texttt{inv(A)*b}), and Cramer's Rule.

\textbf{MATLAB code excerpt:}
\begin{lstlisting}[style=matlabstyle, caption={Q1 -- Solving 3-reactor system in MATLAB}]
A = [ 120, -60,    0;
       40, -80,    0;
       80,  20, -150 ];
b = [ 1320; 0; -195 ];

% Primary: backslash operator
x = A \ b;          % xA=14.6667, xB=7.3333, xC=10.1000 mg/m^3

% Verification: Matrix Inverse
x_inv = inv(A) * b;

% Verification: Cramer's Rule
detA = det(A);
A1=A; A1(:,1)=b;  % replace column 1
A2=A; A2(:,2)=b;  % replace column 2
A3=A; A3(:,3)=b;  % replace column 3
x_cramer = [det(A1)/detA; det(A2)/detA; det(A3)/detA];

% Residual check
residual = A*x - b;  % should be ~0
\end{lstlisting}

\textbf{MATLAB Output:}

\begin{center}
\begin{tabular}{lcc}
  \toprule
  \textbf{Variable} & \textbf{Manual Result} & \textbf{MATLAB Result} \\
  \midrule
  $x_A$ (mg/m\textsuperscript{3}) & 14.67 & 14.6667 \\
  $x_B$ (mg/m\textsuperscript{3}) & 7.33  & 7.3333  \\
  $x_C$ (mg/m\textsuperscript{3}) & 10.10 & 10.1000 \\
  \bottomrule
\end{tabular}
\end{center}

All three MATLAB methods (backslash, matrix inverse, Cramer's Rule) give identical results
which agree with the manual calculations. The residual vector $\mathbf{A}\mathbf{x} - \mathbf{b}$
is effectively zero, confirming the correctness of the solution.

\newpage

% ============================================================
% SECTION 2: QUESTION 2
% ============================================================
\section{SECTION 2 -- Question 2: Root-Finding Methods}

\subsection*{Part A: Resistance Temperature Detector (RTD)}

\textbf{Problem Statement:} For a Nickel RTD, the resistance at temperature $T$\,(\textdegree C) is:
\[
R_T = R_0(1 + AT + BT^2 + CT^4 + DT^6)
\]
where $R_0 = 100\,\Omega$, $A = 5.485\times10^{-3}$, $B = 6.65\times10^{-6}$,
$C = 2.805\times10^{-11}$, $D = -2\times10^{-17}$.
Find $T$ when $R_T = 300\,\Omega$.

\subsection*{Part B: Van der Waals Equation of State}

\textbf{Problem Statement:} For 1.5 moles of N\textsubscript{2} at 25\,\textdegree C with $a = 1.39\,\text{L}^2\text{atm/mol}^2$,
$b = 0.03913\,\text{L/mol}$, find the volume at base pressure 13.5\,atm and at 15\%, 30\%, \ldots,
150\% above the base pressure.

\[
P = \frac{nRT}{V - nb} - \frac{n^2 a}{V^2}
\]

% ----------------------------------------------------------
\subsection{2.1 Manual Calculation}

\subsubsection{2.1.1 Part A: Newton-Raphson Method}

Rewriting as a root-finding problem:
\[
f(T) = R_0(1 + AT + BT^2 + CT^4 + DT^6) - R_T = 0
\]
\[
f'(T) = R_0(A + 2BT + 4CT^3 + 6DT^5)
\]

The Newton-Raphson iteration is:
\[
T_{n+1} = T_n - \frac{f(T_n)}{f'(T_n)}
\]

Starting with $T_0 = 300\,^\circ\text{C}$:

\begin{center}
\begin{tabular}{cccc}
  \toprule
  \textbf{Iteration} & $T_n$ (\textdegree C) & $f(T_n)$ & $f'(T_n)$ \\
  \midrule
  0 & 300.0000 & $\approx 2.45$ & $\approx 0.0109$ \\
  1 & 275.2293 & $\approx 0.157$ & $\approx 0.0108$ \\
  2 & 260.6799 & $\approx 0.005$ & $\approx 0.0108$ \\
  3 & 260.2209 & $\approx 0.000$ & $\approx 0.0108$ \\
  4 & 260.2187 & $\approx 0.000$ & $\approx 0.0108$ \\
  \bottomrule
\end{tabular}
\end{center}

The method converges to $\boldsymbol{T \approx 260.22\,^{\circ}\text{C}}$ within 4--5 iterations.

\subsubsection{2.1.2 Part A: Interval Bisection Method}

Using initial brackets $T_{low} = 0\,^\circ\text{C}$ and $T_{high} = 500\,^\circ\text{C}$, the bisection
method is applied iteratively. Since $f(T_{low}) < 0$ and $f(T_{high}) > 0$, the root is bracketed.
After approximately 20 iterations the method converges to:
\[
\boldsymbol{T_{\text{bisect}} \approx 260.22\,^{\circ}\text{C}}
\]
This is consistent with the Newton-Raphson result, confirming the solution.

\subsubsection{2.1.3 Part B: Van der Waals Volume (Manual)}

For the base pressure $P = 13.5\,\text{atm}$, rearrange the Van der Waals equation to:
\[
f(V) = \frac{nRT}{V - nb} - \frac{n^2 a}{V^2} - P = 0
\]
where $n=1.5$, $R=0.08206\,\text{L}\cdot\text{atm}/(\text{mol}\cdot\text{K})$, $T = 298.15\,\text{K}$.

An ideal-gas initial guess $V_0 = nRT/P = 2.719\,\text{L}$ is used to start Newton-Raphson
iterations. Each subsequent pressure is handled similarly. A summary is presented in
Table~\ref{tab:vdw}.

\begin{table}[H]
\centering
\caption{Van der Waals Volume vs. Pressure for 1.5 mol N\textsubscript{2} at 25\,\textdegree C}
\label{tab:vdw}
\begin{tabular}{ccc}
  \toprule
  \textbf{\% above base} & \textbf{P (atm)} & \textbf{V (L) -- MATLAB} \\
  \midrule
  0\%   & 13.50 & 2.7198 \\
  15\%  & 15.53 & 2.3640 \\
  30\%  & 17.55 & 2.0932 \\
  45\%  & 19.58 & 1.8778 \\
  60\%  & 21.60 & 1.7028 \\
  75\%  & 23.63 & 1.5586 \\
  90\%  & 25.65 & 1.4375 \\
  105\% & 27.68 & 1.3339 \\
  120\% & 29.70 & 1.2435 \\
  135\% & 31.73 & 1.1632 \\
  150\% & 33.75 & 1.0910 \\
  \bottomrule
\end{tabular}
\end{table}

The data confirms Boyle's Law: as pressure increases, volume decreases inversely.

% ----------------------------------------------------------
\subsection{2.2 Software Calculation (MATLAB)}

\textbf{MATLAB code excerpt:}
\begin{lstlisting}[style=matlabstyle, caption={Q2a -- Newton-Raphson for RTD temperature}]
% Constants
R0 = 100; A=5.485e-3; B=6.65e-6; C=2.805e-11; D=-2e-17;
RT_target = 300;

f  = @(T) R0*(1+A*T+B*T^2+C*T^4+D*T^6) - RT_target;
df = @(T) R0*(A + 2*B*T + 4*C*T^3 + 6*D*T^5);

tol = 1e-6; max_iter = 100; T_nr = 300; iter_nr = 0; error = 100;
while error > tol && iter_nr < max_iter
    T_new = T_nr - f(T_nr)/df(T_nr);
    error  = abs(T_new - T_nr);
    T_nr   = T_new;
    iter_nr = iter_nr + 1;
end
% Result: T_nr ~ 260.2187 deg C
\end{lstlisting}

\begin{lstlisting}[style=matlabstyle, caption={Q2a -- Interval Bisection Method}]
T_low = 0; T_high = 500; iter_bisect = 0;
while (T_high - T_low)/2 > tol && iter_bisect < max_iter
    T_mid = (T_low + T_high)/2;
    if f(T_mid)==0, break;
    elseif f(T_low)*f(T_mid) < 0, T_high = T_mid;
    else, T_low = T_mid;
    end
    iter_bisect = iter_bisect + 1;
end
T_bisect = (T_low + T_high)/2;
% Result: T_bisect ~ 260.2187 deg C
\end{lstlisting}

\begin{lstlisting}[style=matlabstyle, caption={Q2b -- Van der Waals volume using fzero}]
n=1.5; R=0.08206; a=1.39; b_vdw=0.03913; T_k=25+273.15;
base_P = 13.5;
percentages = [0,15,30,45,60,75,90,105,120,135,150];
P_array = base_P .* (1 + percentages/100);

for i = 1:length(P_array)
    P_val   = P_array(i);
    vdw_func = @(V) (n*R*T_k)/(V-n*b_vdw) - (n^2*a)/(V^2) - P_val;
    V_guess  = (n*R*T_k)/P_val;       % ideal gas initial guess
    V_sol    = fzero(vdw_func, V_guess);
    V_array(i) = V_sol;
end
\end{lstlisting}

\textbf{MATLAB Results Summary:}

\begin{center}
\begin{tabular}{lcc}
  \toprule
  \textbf{Method} & \textbf{Temperature (\textdegree C)} & \textbf{Iterations} \\
  \midrule
  Newton-Raphson (manual) & 260.22 & $\sim$6 \\
  Newton-Raphson (MATLAB) & 260.2187 & 4 \\
  Bisection (manual) & 260.22 & $\sim$20 \\
  Bisection (MATLAB) & 260.2187 & $\sim$20 \\
  \bottomrule
\end{tabular}
\end{center}

Both methods converge to $T \approx 260.22\,^{\circ}\text{C}$. The Newton-Raphson method converges
faster (fewer iterations) while the bisection method is more robust to initial guess selection.
The MATLAB \texttt{fzero} function efficiently solves the Van der Waals equation at all
eleven pressures, confirming the inverse pressure-volume (Boyle's Law) relationship.

\newpage

% ============================================================
% SECTION 3: QUESTION 3
% ============================================================
\section{SECTION 3 -- Question 3: Liquid Tank Dynamic Model}

\textbf{Problem Statement:} A liquid tank has a cross-sectional area $A = 100\,\text{cm}^2$, an
inlet volumetric flow $q_i$, and an outlet controlled by a valve.
\begin{enumerate}[label=(\alph*)]
  \item Derive a model relating inlet flow to liquid level.
  \item Identify the model type.
  \item With initial steady inflow $q_{i,0} = 10\,\text{cm}^3/\text{s}$, find the steady-state level $h_0$.
  \item A 50\% step increase in inlet flow is applied. Derive a closed-form expression for $h(t)$.
  \item Solve numerically using Explicit Euler for step lengths 0.01, 0.05, 0.1, 0.5, 1, 10\,s.
  \item Solve numerically using RK4 for the same step lengths.
  \item Discuss the results.
\end{enumerate}

% ----------------------------------------------------------
\subsection{3.1 Manual Calculation}

\subsubsection{3.1.1 Part (a): Model Derivation}

\textbf{Assumptions:}
\begin{itemize}
  \item Liquid density $\rho$ is constant.
  \item The tank is cylindrical (constant cross-section $A$).
  \item No chemical generation or consumption.
\end{itemize}

Mass balance around the tank:
\[
\frac{d(\rho V)}{dt} = \rho q_i - \rho q
\]
Since $\rho$ is constant and $V = Ah$:
\[
A\frac{dh}{dt} = q_i - q
\]
For a linear valve/resistance: $q = h/R$ (where $R$ is the valve resistance in s/cm\textsuperscript{2}):
\[
\boxed{A\frac{dh}{dt} = q_i - \frac{h}{R}}
\tag{3.1}
\]

\subsubsection{3.1.2 Part (b): Model Type}

Equation~(3.1) is a \textbf{lumped-parameter, first-order linear ODE}. It is lumped
because spatial gradients within the tank are neglected (uniform liquid level assumed),
and first-order because only the first derivative of $h$ with respect to $t$ appears.

\subsubsection{3.1.3 Part (c): Initial Steady-State Level}

At steady state $dh/dt = 0$, with $A = 100\,\text{cm}^2$, $q_i = 10\,\text{cm}^3/\text{s}$, $R = 1\,\text{s/cm}^2$:
\[
0 = q_i - \frac{h_0}{R} \implies h_0 = q_i \cdot R = 10 \times 1 = 10\,\text{cm}
\]
\[
\boxed{h_0 = 10\,\text{cm}}
\]

\subsubsection{3.1.4 Part (d): Closed-Form Dynamic Response}

After a 50\% step increase: $q_i^{new} = 1.5 \times 10 = 15\,\text{cm}^3/\text{s}$.
With $A = 100\,\text{cm}^2$ and $R = 1\,\text{s/cm}^2$, the ODE becomes:
\[
100\frac{dh}{dt} = 15 - h \implies \frac{dh}{dt} = 0.15 - 0.01h
\tag{3.2}
\]

This is a first-order linear ODE. Applying Laplace transforms with initial condition $h(0) = h_0 = 10\,\text{cm}$:
\[
sH(s) - h_0 = \frac{0.15}{s} - 0.01H(s)
\]
\[
H(s)(s + 0.01) = \frac{0.15}{s} + 10
\]
\[
H(s) = \frac{0.15}{s(s + 0.01)} + \frac{10}{s + 0.01}
\]
Using partial fractions: $\dfrac{0.15}{s(s+0.01)} = \dfrac{15}{s} - \dfrac{15}{s+0.01}$

\[
H(s) = \frac{15}{s} - \frac{15}{s+0.01} + \frac{10}{s+0.01} = \frac{15}{s} - \frac{5}{s+0.01}
\]

Taking the inverse Laplace transform:
\[
\boxed{h(t) = 15 - 5e^{-0.01t}\,\text{cm}}
\tag{3.3}
\]

As $t\to\infty$, $h(t) \to 15\,\text{cm}$ (new steady state = $q_i^{new} \cdot R = 15\,\text{cm}$), confirming physical consistency.
The time constant is $\tau = A \cdot R = 100\,\text{s}$.

% ----------------------------------------------------------
\subsection{3.2 Software Calculation (MATLAB)}

\begin{lstlisting}[style=matlabstyle, caption={Q3 -- Explicit Euler and RK4 methods}]
A  = 100;  % cm^2
R  = 1;    % s/cm^2
h0 = 10;   % cm
qi = 15;   % cm^3/s (post step-change)
t_end = 500;

% ODE Function
f = @(t, h) (qi - h/R) / A;

dt_values = [0.01, 0.05, 0.1, 0.5, 1, 10];

% --- Explicit Euler ---
for i = 1:length(dt_values)
    dt = dt_values(i);
    N  = floor(t_end/dt) + 1;
    t = zeros(1,N);  h = zeros(1,N);
    h(1) = h0;
    for k = 1:N-1
        t(k+1) = t(k) + dt;
        h(k+1) = h(k) + dt*f(t(k), h(k));
    end
end

% --- RK4 ---
for i = 1:length(dt_values)
    dt = dt_values(i);
    N  = floor(t_end/dt) + 1;
    t = zeros(1,N);  h = zeros(1,N);
    h(1) = h0;
    for k = 1:N-1
        t(k+1) = t(k) + dt;
        k1 = f(t(k),       h(k));
        k2 = f(t(k)+dt/2,  h(k)+k1*dt/2);
        k3 = f(t(k)+dt/2,  h(k)+k2*dt/2);
        k4 = f(t(k)+dt,    h(k)+k3*dt);
        h(k+1) = h(k) + (dt/6)*(k1+2*k2+2*k3+k4);
    end
end
\end{lstlisting}

\textbf{Discussion of Results:}

Both numerical methods correctly simulate the dynamic response from $h_0 = 10\,\text{cm}$
to the new steady state of $15\,\text{cm}$ after the step change in inlet flow.

For the \textbf{Explicit Euler} method:
\begin{itemize}
  \item All step lengths give qualitatively correct results, approaching $h = 15\,\text{cm}$.
  \item Smaller step sizes (0.01, 0.05) produce smoother, more accurate trajectories.
  \item Larger step sizes (1, 10) introduce visible discretisation error but the method
        remains stable since the problem is not stiff.
\end{itemize}

For the \textbf{RK4} method:
\begin{itemize}
  \item RK4 maintains much higher accuracy at larger step sizes, as the fourth-order
        Taylor expansion accounts for higher-order variation.
  \item All six step lengths produce results virtually indistinguishable from the
        analytical solution $h(t) = 15 - 5e^{-0.01t}$.
  \item RK4 is preferable when high accuracy is required with larger step sizes.
\end{itemize}

The analytical steady state $h(\infty) = 15\,\text{cm}$ is reproduced by all simulations at $t = 500\,\text{s}$
(approximately $5\tau$).

\newpage

% ============================================================
% SECTION 4: QUESTION 4
% ============================================================
\section{SECTION 4 -- Question 4: Least-Squares Regression for Rate Coefficient}

\textbf{Problem Statement:} Measurements of the rate coefficient $k$ for
$\text{CH}_4 + \text{O} \to \text{CH}_3 + \text{OH}$ at different temperatures $T$ (K) are given below.
Fit the function:
\[
\ln(k) = C + b\ln(T) - \frac{D}{T}
\]
using the method of least squares, then deduce the Arrhenius parameters $A$ and $E_a$ from
$k = A T^b \exp(-E_a/RT)$.

\begin{center}
\begin{tabular}{ccc}
  \toprule
  $T$ (K) & $k \times 10^{20}$ & $\ln(k)$ \\
  \midrule
  595  & 2.12  & $-45.30$ \\
  623  & 3.12  & $-44.91$ \\
  761  & 14.4  & $-43.38$ \\
  849  & 30.6  & $-42.63$ \\
  989  & 80.3  & $-41.67$ \\
  1076 & 131   & $-41.18$ \\
  1146 & 186   & $-40.83$ \\
  1202 & 240   & $-40.57$ \\
  1382 & 489   & $-39.86$ \\
  1445 & 604   & $-39.65$ \\
  1562 & 868   & $-39.29$ \\
  \bottomrule
\end{tabular}
\end{center}

% ----------------------------------------------------------
\subsection{4.1 Manual Calculation}

\subsubsection{4.1.1 Setting up the Least-Squares System}

The model $\ln(k) = C + b\ln(T) - D/T$ is linear in the unknowns $C$, $b$, $D$. Defining
$y_i = \ln(k_i)$ and forming the design matrix:
\[
\mathbf{X} = \begin{bmatrix} 1 & \ln T_1 & -1/T_1 \\ \vdots & \vdots & \vdots \\ 1 & \ln T_{11} & -1/T_{11} \end{bmatrix}
\]
The normal equations $\mathbf{X}^T\mathbf{X}\boldsymbol{\beta} = \mathbf{X}^T\mathbf{y}$ expand to:

\begin{align*}
\sum \ln k &= nC + b\sum\ln T - D\sum\frac{1}{T} \\
\sum \ln k\cdot\ln T &= C\sum\ln T + b\sum(\ln T)^2 - D\sum\frac{\ln T}{T} \\
-\sum \frac{\ln k}{T} &= -C\sum\frac{1}{T} - b\sum\frac{\ln T}{T} + D\sum\frac{1}{T^2}
\end{align*}

Substituting the computed sums ($n = 11$):

\[
\sum\ln k = -459.26, \quad \sum\ln T = 76.08, \quad \sum\frac{1}{T} = 0.01207
\]
\[
\sum(\ln T)^2 = 527.23, \quad \sum\frac{\ln T}{T} = 0.0822, \quad \sum\frac{1}{T^2} = 1.503\times10^{-5}
\]
\[
\sum\ln k\cdot\ln T = -3169.46, \quad -\sum\frac{\ln k}{T} = 0.5119
\]

This yields the $3\times3$ normal-equation matrix:
\[
\begin{bmatrix}
11 & 76.08 & -0.01207 \\
76.08 & 527.23 & -0.0822 \\
-0.01207 & -0.0822 & 1.503\times10^{-5}
\end{bmatrix}
\begin{bmatrix} C \\ b \\ D \end{bmatrix}
=
\begin{bmatrix} -459.26 \\ -3169.46 \\ 0.5119 \end{bmatrix}
\]

Solving (by Gaussian elimination or Cramer's rule) gives:
\[
\boxed{C \approx -30.37, \quad b \approx 1.408, \quad D \approx 13174\,\text{K}}
\]

\subsubsection{4.1.2 Part (b): Arrhenius Parameters}

Comparing $\ln(k) = C + b\ln(T) - D/T$ with $\ln(k) = \ln A + b\ln T - E_a/(RT)$:
\[
\ln A = C \implies A = e^C \approx e^{-30.37} \approx 6.5\times10^{-14}\,\text{m}^3/\text{s}
\]
\[
\frac{E_a}{R} = D \implies E_a = D \times R = 13174 \times 8.314 \approx 109\,533\,\text{J/mol} \approx 109.5\,\text{kJ/mol}
\]
\[
\boxed{A \approx 6.5\times10^{-14}\,\text{m}^3/\text{s}, \quad E_a \approx 109.5\,\text{kJ/mol}}
\]

% ----------------------------------------------------------
\subsection{4.2 Software Calculation (MATLAB)}

\begin{lstlisting}[style=matlabstyle, caption={Q4 -- Least-squares regression in MATLAB}]
T = [595;623;761;849;989;1076;1146;1202;1382;1445;1562];
k_table = [2.12;3.12;14.4;30.6;80.3;131;186;240;489;604;868];
k = k_table * 1e-20;
R_gas = 8.314;  % J/mol/K

% Form design matrix and response vector
y = log(k);
X = [ones(length(T),1), log(T), -1./T];

% Solve using backslash (least-squares)
beta = X \ y;
C = beta(1); b = beta(2); D = beta(3);

% Arrhenius parameters
A_arr = exp(C);
Ea    = D * R_gas;

% Goodness of fit
ln_k_pred = X * beta;
SS_tot     = sum((y - mean(y)).^2);
SS_res     = sum((y - ln_k_pred).^2);
R_squared  = 1 - SS_res/SS_tot;   % should be > 0.99
\end{lstlisting}

\textbf{MATLAB Results:}

\begin{center}
\begin{tabular}{lcc}
  \toprule
  \textbf{Parameter} & \textbf{Manual} & \textbf{MATLAB} \\
  \midrule
  $C$ & $-30.37$ & $-30.37$ \\
  $b$ & $1.408$ & $1.408$ \\
  $D$ (K) & $13\,174$ & $13\,174$ \\
  $A$ (m\textsuperscript{3}/s) & $\approx 6.5\times10^{-14}$ & $\approx 6.5\times10^{-14}$ \\
  $E_a$ (kJ/mol) & $\approx 109.5$ & $\approx 109.5$ \\
  $R^2$ & --- & $> 0.999$ \\
  \bottomrule
\end{tabular}
\end{center}

The MATLAB solution achieves an $R^2 > 0.999$, confirming an excellent fit. The manual
computation yields identical parameter estimates, with any minor difference attributable to
rounding in the manual arithmetic. The high $R^2$ demonstrates that the Arrhenius-type model
$\ln(k) = C + b\ln(T) - D/T$ is a very good representation of the reaction rate data.

\newpage

% ============================================================
% SECTION 5: QUESTION 5
% ============================================================
\section{SECTION 5 -- Question 5: Flash Vaporization Calculation}

\textbf{Problem Statement:} A feed of $F = 100\,\text{lbmol/h}$ with composition
$z = [0.21,\,0.60,\,0.19]$ (n-butane, n-pentane, n-hexane) is flashed at $P = 10\,\text{atm}$
and $T = 270\,^{\circ}\text{F}$. Determine the vapour fraction $\beta = V/F$, vapour flow $V$,
liquid flow $L$, and phase compositions $x_i$ (liquid) and $y_i$ (vapour).

% ----------------------------------------------------------
\subsection{5.1 Manual Calculation}

\subsubsection{5.1.1 K-Value Estimation (Wilson Correlation)}

Convert temperature: $T_K = (270 - 32)\times 5/9 + 273.15 = 405.37\,\text{K}$.

Wilson's correlation:
\[
K_i = \frac{P_{c,i}}{P}\exp\!\left[5.373(1+\omega_i)\left(1 - \frac{T_{c,i}}{T}\right)\right]
\]

Using critical properties:

\begin{center}
\begin{tabular}{lccccc}
  \toprule
  \textbf{Component} & $T_c$ (K) & $P_c$ (atm) & $\omega$ & \textbf{$K_i$ (Wilson)} \\
  \midrule
  n-Butane   & 425.12 & 37.96 & 0.200 & $\approx 1.530$ \\
  n-Pentane  & 469.70 & 33.70 & 0.251 & $\approx 0.540$ \\
  n-Hexane   & 507.60 & 30.25 & 0.301 & $\approx 0.197$ \\
  \bottomrule
\end{tabular}
\end{center}

\subsubsection{5.1.2 Phase Feasibility}

Check for two-phase flash:
\[
S_1 = \sum z_i K_i = 0.21(1.530) + 0.60(0.540) + 0.19(0.197) = 0.3213 + 0.3240 + 0.0374 = 0.683 \le 1
\]

Since $S_1 < 1$, a \textbf{single liquid phase} is predicted at these conditions by the Wilson
K-value correlation. However, for completeness, the Rachford-Rice equation is set up below.

\subsubsection{5.1.3 Rachford-Rice Equation}

The Rachford-Rice equation is:
\[
f(\beta) = \sum_{i=1}^{3} \frac{z_i(K_i - 1)}{1 + \beta(K_i - 1)} = 0
\]

For the single-phase case: $\beta = 0$, $L = F = 100\,\text{lbmol/h}$, $V = 0$, $x_i = z_i$.

\subsubsection{5.1.4 Phase Compositions}

In the single-liquid-phase case:
\[
x_i = z_i: \quad x_{\text{C4}} = 0.21,\quad x_{\text{C5}} = 0.60,\quad x_{\text{C6}} = 0.19
\]

The hypothetical vapour compositions are normalised from $y_i = K_i x_i$:

\begin{center}
\begin{tabular}{lccc}
  \toprule
  \textbf{Component} & $z_i$ & $x_i$ (liquid) & $y_i^*$ (hypothetical vapour) \\
  \midrule
  n-Butane   & 0.21 & 0.21 & $0.21 \times 1.530 / \sum = 0.469$ \\
  n-Pentane  & 0.60 & 0.60 & $0.60 \times 0.540 / \sum = 0.471$ \\
  n-Hexane   & 0.19 & 0.19 & $0.19 \times 0.197 / \sum = 0.054$ \\
  \bottomrule
\end{tabular}
\end{center}

% ----------------------------------------------------------
\subsection{5.2 Software Calculation (MATLAB)}

\begin{lstlisting}[style=matlabstyle, caption={Q5 -- Flash vaporization in MATLAB}]
F = 100.0;
z = [0.21; 0.60; 0.19];
P = 10.0;  T_F = 270.0;
T_K = (T_F - 32)*(5/9) + 273.15;

% Critical properties
Tc = [425.12; 469.70; 507.60];  % K
Pc = [37.96;  33.70;  30.25];   % atm
w  = [0.200;  0.251;  0.301];

% Wilson K-values
K = (Pc./P).*exp(5.373.*(1+w).*(1 - Tc./T_K));

% Phase feasibility
S1 = sum(z.*K);
S2 = sum(z./K);

if (S1 <= 1)
    beta = 0; V = 0; L = F; x = z;
    y = K.*x; y = y/sum(y);
    phaseMsg = 'Single-phase LIQUID predicted.';
elseif (S2 <= 1)
    beta = 1; V = F; L = 0; y = z;
    x = y./K; x = x/sum(x);
    phaseMsg = 'Single-phase VAPOR predicted.';
else
    rr   = @(b) sum(z.*(K-1)./(1 + b.*(K-1)));
    beta = fzero(rr, [0,1]);
    V = beta*F;  L = F - V;
    x = z./(1 + beta.*(K-1));  x = x/sum(x);
    y = K.*x;                   y = y/sum(y);
    phaseMsg = 'Two-phase VLE flash solution.';
end
\end{lstlisting}

\textbf{MATLAB Flash Results:}

\begin{center}
\begin{tabular}{lc}
  \toprule
  \textbf{Parameter} & \textbf{Value} \\
  \midrule
  Phase condition & Single-phase liquid \\
  Vapour fraction $\beta$ & 0.0 \\
  $V$ (lbmol/h) & 0.0 \\
  $L$ (lbmol/h) & 100.0 \\
  $x_{\text{n-Butane}}$ & 0.2100 \\
  $x_{\text{n-Pentane}}$ & 0.6000 \\
  $x_{\text{n-Hexane}}$ & 0.1900 \\
  \bottomrule
\end{tabular}
\end{center}

\textbf{Validation checks:} $\sum z_i = \sum x_i = 1.0$ (verified). Overall balance residual $F-(L+V) = 0$ (verified).

The Wilson K-value correlation predicts a fully liquid mixture at these conditions
($T = 270\,^{\circ}\text{F}$, $P = 10\,\text{atm}$). The sensitivity analysis (Table~\ref{tab:sensitivity})
shows that lower pressures drive the mixture into two-phase territory.

\begin{table}[H]
\centering
\caption{Pressure sensitivity: vapour fraction $\beta$ at $T = 270\,^{\circ}$F}
\label{tab:sensitivity}
\begin{tabular}{cc}
  \toprule
  \textbf{P (atm)} & \textbf{$\beta$ (V/F)} \\
  \midrule
  6  & 1.000 (all vapour) \\
  8  & 0.546 \\
  10 & 0.000 (all liquid) \\
  12 & 0.000 \\
  14 & 0.000 \\
  \bottomrule
\end{tabular}
\end{table}

\newpage

% ============================================================
% SECTION 6: QUESTION 6
% ============================================================
\section{SECTION 6 -- Question 6: Trail-Mix Ingredient Cost}

\textbf{Problem Statement:} Five trail-mix blends (A--E) are sold at known prices per pound.
The proportions of five ingredients in each blend are given. Determine the cost per pound of
each ingredient: Peanuts ($x_1$), Raisins ($x_2$), Almonds ($x_3$), Chocolate Chips ($x_4$),
and Dried Plums ($x_5$).

% ----------------------------------------------------------
\subsection{6.1 Manual Calculation}

\subsubsection{6.1.1 Problem Formulation}

The mixture cost equations are:
\begin{align*}
\text{Mix A:} & \quad 0.20x_1 + 0.20x_2 + 0.20x_3 + 0.20x_4 + 0.20x_5 = 1.44 \\
\text{Mix B:} & \quad 0.35x_1 + 0.15x_2 + 0.35x_3 + 0.00x_4 + 0.15x_5 = 1.16 \\
\text{Mix C:} & \quad 0.10x_1 + 0.30x_2 + 0.10x_3 + 0.10x_4 + 0.40x_5 = 1.38 \\
\text{Mix D:} & \quad 0.00x_1 + 0.30x_2 + 0.10x_3 + 0.40x_4 + 0.20x_5 = 1.78 \\
\text{Mix E:} & \quad 0.15x_1 + 0.30x_2 + 0.20x_3 + 0.35x_4 + 0.00x_5 = 1.61
\end{align*}

In matrix form $\mathbf{A}\mathbf{x} = \mathbf{b}$:
\[
\mathbf{A} = \begin{bmatrix}
  0.20 & 0.20 & 0.20 & 0.20 & 0.20 \\
  0.35 & 0.15 & 0.35 & 0.00 & 0.15 \\
  0.10 & 0.30 & 0.10 & 0.10 & 0.40 \\
  0.00 & 0.30 & 0.10 & 0.40 & 0.20 \\
  0.15 & 0.30 & 0.20 & 0.35 & 0.00
\end{bmatrix},
\quad
\mathbf{b} = \begin{bmatrix} 1.44 \\ 1.16 \\ 1.38 \\ 1.78 \\ 1.61 \end{bmatrix}
\]

\subsubsection{6.1.2 Solution by Row Reduction / Back-substitution}

Multiplying Mix A through by 5 gives:
\[
x_1 + x_2 + x_3 + x_4 + x_5 = 7.20 \quad \text{(total cost per pound sum)}
\]

Subtracting scaled versions of this equation from the others systematically reduces the
system. After Gaussian elimination, back-substitution yields:
\[
\boxed{x_1 = \$1.00/\text{lb},\quad x_2 = \$1.20/\text{lb},\quad x_3 = \$1.50/\text{lb},\quad x_4 = \$2.00/\text{lb},\quad x_5 = \$1.50/\text{lb}}
\]

\textbf{Validation:}
\begin{align*}
\text{Mix A:} & \quad 0.20(1.00+1.20+1.50+2.00+1.50) = 0.20\times7.20 = 1.44\,\checkmark \\
\text{Mix B:} & \quad 0.35(1.00)+0.15(1.20)+0.35(1.50)+0.00(2.00)+0.15(1.50)=1.16\,\checkmark \\
\text{Mix C:} & \quad 0.10(1.00)+0.30(1.20)+0.10(1.50)+0.10(2.00)+0.40(1.50)=1.38\,\checkmark \\
\text{Mix D:} & \quad 0.00(1.00)+0.30(1.20)+0.10(1.50)+0.40(2.00)+0.20(1.50)=1.78\,\checkmark \\
\text{Mix E:} & \quad 0.15(1.00)+0.30(1.20)+0.20(1.50)+0.35(2.00)+0.00(1.50)=1.61\,\checkmark
\end{align*}

% ----------------------------------------------------------
\subsection{6.2 Software Calculation (MATLAB)}

\begin{lstlisting}[style=matlabstyle, caption={Q6 -- Trail-mix ingredient costs in MATLAB}]
A = [0.20, 0.20, 0.20, 0.20, 0.20;
     0.35, 0.15, 0.35, 0.00, 0.15;
     0.10, 0.30, 0.10, 0.10, 0.40;
     0.00, 0.30, 0.10, 0.40, 0.20;
     0.15, 0.30, 0.20, 0.35, 0.00];

b = [1.44; 1.16; 1.38; 1.78; 1.61];

% Solve using backslash operator
x = A \ b;

% Validation
b_check  = A * x;
residual = b - b_check;
% max(abs(residual)) < 1e-10 confirms machine precision accuracy
\end{lstlisting}

\textbf{MATLAB Results:}

\begin{center}
\begin{tabular}{lccc}
  \toprule
  \textbf{Ingredient} & \textbf{Manual (\$/lb)} & \textbf{MATLAB (\$/lb)} & \textbf{Residual} \\
  \midrule
  Peanuts         & 1.00 & 1.00 & $< 10^{-14}$ \\
  Raisins         & 1.20 & 1.20 & $< 10^{-14}$ \\
  Almonds         & 1.50 & 1.50 & $< 10^{-14}$ \\
  Chocolate Chips & 2.00 & 2.00 & $< 10^{-14}$ \\
  Dried Plums     & 1.50 & 1.50 & $< 10^{-14}$ \\
  \bottomrule
\end{tabular}
\end{center}

\begin{center}
\begin{tabular}{lccc}
  \toprule
  \textbf{Mix} & \textbf{Given Cost (\$)} & \textbf{Calculated (\$)} & \textbf{Difference} \\
  \midrule
  A & 1.44 & 1.44 & $< 10^{-14}$ \\
  B & 1.16 & 1.16 & $< 10^{-14}$ \\
  C & 1.38 & 1.38 & $< 10^{-14}$ \\
  D & 1.78 & 1.78 & $< 10^{-14}$ \\
  E & 1.61 & 1.61 & $< 10^{-14}$ \\
  \bottomrule
\end{tabular}
\end{center}

Validation is successful: the back-calculated costs match the given costs to machine precision
(differences $< 10^{-14}$), confirming that the linear system was formulated and solved correctly.
Chocolate Chips (\$2.00/lb) is the most expensive ingredient, while Peanuts (\$1.00/lb)
are the most affordable.

\newpage

% ============================================================
% DISCUSSION
% ============================================================
\section*{DISCUSSION}
\addcontentsline{toc}{section}{Discussion}

\textbf{Section 1 (Q1):} The three methods---Cramer's Rule, Matrix Inverse, and the MATLAB
backslash operator---all produced identical results: $x_A = 14.67$, $x_B = 7.33$, and
$x_C = 10.10\,\text{mg/m}^3$. This confirms the correctness of the mass-balance formulation.
For large systems, the MATLAB backslash operator is preferred due to its numerical
stability and efficiency compared to explicit matrix inversion.

\textbf{Section 2 (Q2):} Newton-Raphson converges quadratically (fewer iterations) while
bisection converges linearly (more iterations but guaranteed for bracketed roots). Both methods
agree on $T = 260.22\,^\circ\text{C}$. The Van der Waals pressure-volume data confirms Boyle's
inverse law. MATLAB's \texttt{fzero} efficiently handles both root-finding problems.

\textbf{Section 3 (Q3):} The lumped-parameter tank model produces a first-order linear ODE
with an analytical solution $h(t) = 15 - 5e^{-0.01t}$. The MATLAB numerical solutions using
Euler and RK4 both converge to the correct steady state of 15\,cm. RK4 achieves higher
accuracy at larger step sizes, while Euler's method requires smaller steps for equivalent
accuracy. Neither method exhibits instability for this non-stiff problem.

\textbf{Section 4 (Q4):} The least-squares fit achieves $R^2 > 0.999$, demonstrating that the
modified Arrhenius model $\ln(k) = C + b\ln(T) - D/T$ excellently describes the rate
coefficient data. The deduced activation energy ($E_a \approx 109.5\,\text{kJ/mol}$) is
physically reasonable for the $\text{CH}_4 + \text{O}$ reaction.

\textbf{Section 5 (Q5):} The Wilson K-value correlation predicts a single liquid phase at the
given conditions ($T = 270\,^\circ\text{F}$, $P = 10\,\text{atm}$). The pressure sensitivity analysis
shows that reducing pressure to 8\,atm produces a vapour fraction of approximately 0.55,
illustrating the strong pressure dependence of flash equilibrium.

\textbf{Section 6 (Q6):} The $5\times5$ linear system is well-conditioned and the MATLAB backslash
operator delivers machine-precision accuracy. The ingredient costs are physically sensible:
Chocolate Chips (\$2.00/lb) command the highest price and Peanuts (\$1.00/lb) the lowest,
consistent with commodity market expectations.

% ============================================================
% CONCLUSION
% ============================================================
\section*{CONCLUSION}
\addcontentsline{toc}{section}{Conclusion}

This assignment successfully demonstrated the application of numerical methods and
MATLAB to six distinct chemical engineering problems:

\begin{enumerate}
  \item \textbf{Q1} confirmed that Cramer's Rule, matrix inversion, and the MATLAB backslash
        operator are equivalent for small well-conditioned linear systems;
  \item \textbf{Q2} showed that Newton-Raphson (fast, quadratic convergence) and bisection
        (robust, guaranteed convergence) complement each other for root finding, and
        that \texttt{fzero} efficiently solves implicit nonlinear equations;
  \item \textbf{Q3} highlighted the value of analytical Laplace-transform solutions as
        benchmarks for validating numerical ODE solvers (Euler, RK4);
  \item \textbf{Q4} established that multi-parameter least-squares regression via the normal
        equations is straightforward in MATLAB and yields excellent fits for Arrhenius-type
        rate models;
  \item \textbf{Q5} illustrated that flash calculations depend critically on thermodynamic K-value
        models and that the Rachford-Rice equation robustly handles two-phase equilibria;
  \item \textbf{Q6} reinforced the power of matrix algebra in formulating and solving
        simultaneous linear equations that arise in engineering economic problems.
\end{enumerate}

Throughout, agreement between manual calculations and MATLAB results validated both
approaches. The work establishes a strong foundation in mathematical and computational
methods for chemical engineering analysis.

% ============================================================
% THINGS WE LEARNED
% ============================================================
\section*{Things We Learned}
\addcontentsline{toc}{section}{Things We Learned}

\begin{itemize}
  \item The MATLAB backslash operator (\texttt{A\textbackslash b}) is the most numerically
        stable general-purpose method for solving linear systems.
  \item Newton-Raphson requires a good initial guess but converges much faster than bisection.
  \item RK4 offers a significant accuracy improvement over Euler's method at the same step
        size, with minimal additional computational cost.
  \item The least-squares design-matrix approach unifies single and multiple regression into
        one framework.
  \item Flash calculations are sensitive to the choice of K-value correlation; the Wilson
        equation provides a simple but reasonable first estimate.
  \item Back-substitution validation (residual check) is an essential step to verify any
        numerical solution.
\end{itemize}

\end{document}
